% Part: second-order-logic
% Chapter: syntax-and-semantics
% Section: terms-formulas

\documentclass[../../../include/open-logic-section]{subfiles}

\begin{document}

\olfileid{sol}{syn}{frm}

\olsection{الحدود و\printtoken{P}{formula}}

كما في منطق الرتبة الأولى، تبنى تعبيرات منطق الرتبة الثانية من معجم
أساسي يحتوي \emph{!!{variable}s} و\emph{!!{constant}s}
و\emph{!!{predicate}s}، وأحيانًا \emph{!!{function}s}. ومنها، مع
الروابط المنطقية والمكممات وعلامات الترقيم كالأقواس والفواصل، تتكون
\emph{الحدود} و\emph{!!{formula}s}. والفرق هو أن منطق الرتبة الثانية
يحتوي، فضلًا عن متغيرات الأفراد، متغيرات للعلاقات والدوال، ويسمح
بالتكميم عليها. ولذلك تكون الرموز المنطقية لمنطق الرتبة الثانية هي رموز
منطق الرتبة الأولى مضافًا إليها ما يأتي:

\begin{enumerate}
\item مجموعة !!{denumerable} من !!{variable}s العلاقات من الرتبة
  الثانية لكل رتبة~$n$: $\Obj V_0^n$، و$\Obj V_1^n$، و$\Obj V_2^n$،
  و\dots
\item مجموعة !!{denumerable} من !!{variable}s الدوال من الرتبة
  الثانية: $\Obj u_0^n$، و$\Obj u_1^n$، و$\Obj u_2^n$، و\dots
\end{enumerate}

وكما نستعمل~$x$ و$y$ و$z$ متغيرات فوقية لمتغيرات الرتبة الأولى
$\Obj v_i$، سنستعمل~$X$ و$Y$ و$Z$، إلخ، متغيرات فوقية لـ
$\Obj V_i^n$، ونستعمل~$u$ و$v$، إلخ، متغيرات فوقية لـ$\Obj u_i^n$.

\begin{explain}
تحدد الرموز غير المنطقية للغة من الرتبة الثانية بالطريقة نفسها التي
تحدد بها لغة من الرتبة الأولى: بسرد !!{constant}s و!!{function}s
و!!{predicate}s فيها.

تدرج الـ!!{identity}~$\eq$ عادةً في منطق الرتبة الأولى. وفي منطق
الرتبة الأولى تكون الرموز غير المنطقية للغة~$\Lang{L}$ حاسمة لكي
نستطيع التعبير عن أي شيء مهم. وثمة بالطبع !!{sentence}s لا تستعمل
رموزًا غير منطقية، لكن يصعب قول شيء مهم باستعمال~$\eq$ وحدها. أما في
منطق الرتبة الثانية، فلأن لدينا موردًا غير محدود من متغيرات العلاقات
والدوال، نستطيع أن نقول كل ما يمكن قوله في لغة من الرتبة الأولى حتى
من دون مورد خاص من الرموز غير المنطقية.
\end{explain}

\begin{defn}[حدود الرتبة الثانية]
تعرف مجموعة \emph{حدود الرتبة الثانية} للغة~$\Lang L$، أي
$\TrmSOL[L]$، بإضافة البند الآتي إلى
\olref[fol][syn][frm]{defn:terms}:
\begin{enumerate}
\item إذا كان~$u$ متغيرًا داليًا ذا~$n$ مواضع، وكانت
$t_1$، و\dots، و$t_n$ حدودًا، فإن~$\Atom{u}{t_1, \ldots, t_n}$ حد.
\end{enumerate}
\end{defn}

\begin{explain}
وهكذا يبدو حد الرتبة الثانية كحد الرتبة الأولى تمامًا، غير أنه حيث
يحتوي حد الرتبة الأولى على !!a{function}~$\Obj{f^n_i}$ قد يحتوي حد
الرتبة الثانية بدلًا منها على متغير دالي~$\Obj{u^n_i}$.
\end{explain}

\begin{defn}[\usetoken{s}{formula} الرتبة الثانية]
تعرف مجموعة \emph{!!{formula}s الرتبة الثانية}~$\FrmSOL[L]$ للغة
$\Lang L$ بإضافة البنود الآتية إلى
\olref[fol][syn][frm]{defn:formulas}:
\begin{enumerate}
\item إذا كان~$X$ متغيرًا محموليًا ذا~$n$ مواضع، وكانت
$t_1$، و\dots، و$t_n$ حدودًا من الرتبة الثانية للغة~$\Lang L$، فإن
$\Atom{X}{t_1,\ldots, t_n}$ !!{formula} ذرية.

\tagitem{prvAll}{إذا كانت~$!A$ !!a{formula} وكان~$u$ متغيرًا داليًا،
  فإن~$\lforall[u][!A]$ !!a{formula}.}{}

\tagitem{prvAll}{إذا كانت~$!A$ !!a{formula} وكان~$X$ متغيرًا محموليًا،
  فإن~$\lforall[X][!A]$ !!a{formula}.}{}

\tagitem{prvEx}{إذا كانت~$!A$ !!a{formula} وكان~$u$ متغيرًا داليًا،
  فإن~$\lexists[u][!A]$ !!a{formula}.}{}

\tagitem{prvEx}{إذا كانت~$!A$ !!a{formula} وكان~$X$ متغيرًا محموليًا،
  فإن~$\lexists[X][!A]$ !!a{formula}.}{}
\end{enumerate}
\end{defn}

\end{document}
